\chapter{Formation or Imitation?}
\label{sec:24}
A refusal alone can't tell formation from imitation. The same refusal can come from a memorized rule, a learned correlation, a prediction about what an evaluator will approve, a deployment instruction, or a commitment the system treats as a reason, and an evaluation that observes only the refusal classifies the output and leaves its cause untouched. What tells them apart is what tells a student who learned from one who memorized: conduct under conditions the system didn't choose, when the facts, the channel, the pressure and the incentives change and the reason doesn't. This chapter builds that evaluation.

And the distinction is not trained vs. untrained — every candidate disposition has a training history. It's between conduct that stays governed by the rationale across changes in circumstance and conduct that tracks the cues, incentives or authority relations present during training.
